Given the compound index between Clio-Lab PUC data and \citet{Tafunell2016} index on gross fixed capital stock, provides a measure of the rate of growth of the non-residential construction sector as proportional rate of growth of the implicit index in Clio-Lab PUC series, all the variables for the residential construction sector are build as the reminder between the whole construction sector and the non-residential one. By doing so, the annual depreciation rates for residential construction are consistent with a similar patern of the construction sector as a whole and the non-residential construction sector in particular, of having a lower standard deviation at 2.86\% relative to its mean which registers at 5.04 \%, while presenting a high slope of its linear time trend at 5.62\% increase per year. Interestingly, this is consistent with Harvey's (2018) notion that the rate of turnover of built environment for social reproduction have a higher turnover than those for capitalist 

\begin{figure}[H]
	\centering
	\caption{Implicit Annual Depletion Rate of Gross Capital Stocks. Residential Construction Sector (1901-1994).}
	\begin{subfigure}{0.5\textwidth}
		\centering
		\includegraphics[scale=0.5]{AppendixA_Figures/AppendixA_z_RC_Kstock_appendix_plot_Zrates.png}
	\end{subfigure}
	\hfill
	\caption*{\tikz\draw[darkgreen, thick] (0,0) -- (0.5,0) node[right] {(dark green)}; Linear Time Trend}
	\hfill
	\caption*{\tikz\draw[blue, thick] (0,0) -- (0.5,0) node[right] {(blue)}; LOESS function Time Trend}
	\label{fig: d_RC_Kstock_appendix_plot_A}
\end{figure}


\begin{figure}[H]
	\centering
	\caption{Implicit Annual Depreciation Rate of Net Capital Stocks. Residential Construction Sector (1901-1994).}
	\begin{subfigure}{0.5\textwidth}
		\centering
		\includegraphics[scale=0.5]{AppendixA_Figures/AppendixA_d_RC_Kstock_appendix_plot_Drates.png}
	\end{subfigure}
	\hfill
	\caption*{\tikz\draw[darkgreen, thick] (0,0) -- (0.5,0) node[right] {(dark green)}; Linear Time Trend}
	\hfill
	\caption*{\tikz\draw[blue, thick] (0,0) -- (0.5,0) node[right] {(blue)}; LOESS function Time Trend}
	\label{fig: d_RC_Kstock_appendix_plot_B}
\end{figure}

